\lecture{5.7}{Practical Interpretation of Super Learner Results}
\label{lec:5-7}

\subsection{Candidate Algorithms}

A Super Learner analysis should report the candidate library. The library is part of
the estimator. Two analyses with different libraries are using different statistical
procedures, even if both are called Super Learner.

Common candidate algorithms include:
\begin{itemize}
\item generalized linear models;
\item models with interactions or polynomial terms;
\item penalized regressions such as lasso and ridge;
\item generalized additive models;
\item multivariate adaptive regression splines;
\item regression trees and random forests;
\item gradient boosting;
\item Bayesian additive regression trees;
\item support vector machines;
\item neural networks;
\item simple empirical or intercept-only benchmarks.
\end{itemize}

The library should include algorithms with different levels of flexibility. It is often
useful to include simple benchmark models because they can perform well when the
sample is small or the signal is simple.

\subsection{Cross-Validated Risk Tables}

Super Learner output often includes a table of cross-validated risks. For squared-error
loss, lower risk means lower estimated out-of-sample mean squared error. For negative
log-likelihood loss, lower risk means better predicted probabilities.

A risk table might have the following conceptual form:
\[
\begin{array}{lcc}
\toprule
\text{Algorithm} & \text{CV risk for }\Qbar & \text{CV risk for }g\\
\midrule
\text{GLM} & 1.000 & 1.000\\
\text{GAM} & 0.994 & 0.987\\
\text{Random forest} & 0.997 & 0.989\\
\text{BART} & 0.988 & 4.462\\
\text{Super Learner} & 0.986 & 0.983\\
\bottomrule
\end{array}
\]
Risks are sometimes reported relative to a reference method, such as GLM. A relative
risk below 1 means the algorithm had lower cross-validated risk than the reference.
A relative risk above 1 means it performed worse by the chosen loss.

\subsection{Super Learner Weights}

The ensemble weights describe how the final prediction function is assembled:
\[
\Qbar_{n,SL}(a,w)=\sum_{k=1}^K\omega_k\Qbar_{n,k}(a,w).
\]
A large weight means that an algorithm contributes substantially to the final ensemble
prediction. A zero weight means that the algorithm did not improve the ensemble risk
given the other algorithms in the library.

Weights should not be interpreted as posterior probabilities that a model is true. A
Super Learner library is not a list of mutually exclusive scientific models. It is a
collection of prediction algorithms. The weights are chosen for predictive performance
under a loss function, not for causal interpretation.

\subsection{Why Zero Weights Are Common}

Zero weights are common because many algorithms make similar predictions. If two
algorithms are highly correlated, the ensemble may need only one of them. An
algorithm can also receive zero weight because it predicts poorly, because it is
redundant with better algorithms, or because adding it would increase validation risk.

Zero weight does not necessarily mean an algorithm is useless in all datasets. It means
that, in this dataset and this library, the cross-validated metalearner did not find it
helpful for reducing the chosen risk.

\begin{keyidea}
Super Learner weights are performance weights, not truth weights. They summarize
how candidate predictions were combined to minimize cross-validated loss.
\end{keyidea}

\subsection{Outcome Regression vs Propensity Score Performance}

The same algorithm can perform differently for $\Qbar$ and $g$. The outcome
regression predicts $Y$ from $(A,W)$, while the propensity score predicts $A$ from
$W$. These are different statistical tasks.

For example, a flexible tree-based method might capture nonlinear outcome patterns
well but produce unstable propensity scores near zero or one. A generalized additive
model might predict treatment assignment well because the treatment mechanism is
smooth, while offering little improvement for a noisy outcome.

This distinction is important for causal estimators. In g-computation, the outcome
regression is the central nuisance function. In IPTW, the propensity score is central.
In doubly robust estimators introduced later, both nuisance functions matter, and their
estimation quality can affect bias, variance, and confidence interval behavior.

\subsection{Reporting Model Libraries and Sensitivity Analyses}

Transparent reporting should include:
\begin{itemize}
\item the candidate algorithms;
\item important tuning parameter grids;
\item the loss function;
\item the number of cross-validation folds;
\item whether candidates used screening or preprocessing;
\item the cross-validated risks;
\item the selected discrete learner or ensemble weights;
\item any convergence failures or algorithms removed before analysis.
\end{itemize}

Sensitivity analyses are useful because Super Learner performance depends on the
library. A narrow library cannot choose a method that was not included. Analysts
should examine whether conclusions change when the library is expanded, when
highly unstable algorithms are removed, or when alternative reasonable losses are
used.

\begin{implementationbox}
Practical reporting template:
\begin{lstlisting}
Outcome regression:
  Loss: squared-error or negative log-likelihood
  Folds: V-fold cross-validation
  Library: GLM, GAM, lasso, random forest, ...
  Selected/weighted algorithms: report CV risks and weights

Propensity score:
  Loss: Bernoulli negative log-likelihood
  Folds: V-fold cross-validation
  Library: GLM, GAM, lasso, random forest, ...
  Diagnostics: predicted probabilities and weight summaries
\end{lstlisting}
\end{implementationbox}

\begin{estimationbox}
When Super Learner is used inside a causal analysis, it should be described as part
of the estimator of the causal parameter. Report the causal assumptions separately
from the regression performance. Good cross-validated prediction does not verify
exchangeability, consistency, or positivity.
\end{estimationbox}
